% RESEARCH_DRAFT_SHA256: 35742cff0725c761967e87c64145476cb3a1166abda72587fcf3a239e7d93e3a
% Layout and source-reference wrappers expanded for the web renderer.
\documentclass[11pt]{article}
\usepackage[margin=25mm]{geometry}
\usepackage{fontspec,unicode-math}
\setmainfont{Latin Modern Roman}
\setmathfont{Latin Modern Math}
\usepackage{amsmath,amssymb,mathtools}
\usepackage{xurl,hyperref}
\hypersetup{colorlinks=true,urlcolor=blue,linkcolor=blue}
\setcounter{secnumdepth}{0}
\setlength{\parindent}{0pt}
\setlength{\parskip}{5pt}
\setlength{\emergencystretch}{3em}
\newcommand{\sref}[2]{\hyperlink{source:#1:#2}{[S#2]}}
\newcommand{\cataloguescope}[1]{\paragraph{Source question.}#1}
\newcommand{\floor}[1]{\left\lfloor#1\right\rfloor}
\begin{document}
\section*{Research-attempt provenance}
This standalone research attempt retains the complete original definition body below,
including its historical source and preparation notes. Those notes describe the
earlier catalogue preparation; the new attempt and its verification are identified
separately after that body. No original ID, catalogue status, or source statement is changed.

% UnsolvedMath ID 30001568; source catalogue code OWR-4426-005
\section{Symmetry-Preserving Evaluation of Orbitwise Generating Functions}\label{30001568}
{\small UnsolvedMath ID 30001568 · Computer Science}
\subsection{Definitions and mathematical statement}

\paragraph{Shared notation.}Standard number systems are denoted by $\mathbb N,\mathbb Z,\mathbb Q,\mathbb R,\mathbb C$; index ranges, computational models and specialized notation follow the cited problem source and the definitions below.


For integer points $a$ in a polyhedron, a generating function sums monomials $z^a$. Orbitwise organization groups terms according to the symmetry action. Rational expressions may have apparent poles at $z=(1,\ldots,1)$ even when a bounded-set count has a finite value.

A polyhedron is an intersection of finitely many closed affine half-spaces; a polytope is a bounded polyhedron. Its faces are its intersections with supporting hyperplanes.
\paragraph{Statement recorded in the supplied source.}
Given an orbitwise rational generating function for the integer points of a symmetric polyhedron, how can its value at the singular point $z=(1,\ldots,1)$ be evaluated in an orbitwise way while preserving the symmetry? \href{https://doi.org/10.4171/owr/2010/38}{[S2]}
\paragraph{Scope qualification.}
Boundedness or a regularization convention is needed when the actual lattice-point set is infinite; the supplied question does not specify this.
\paragraph{Source status.}
The supplied catalogue labels this record partially solved. Its source-status metadata was checked against the supplied snapshot on 2026-10-11. This import does not independently establish current open status or mathematical correctness. The catalogue literature assessment is dated 2026-08-21.
\subsection{Short English statement}
Evaluate the symmetry-organized generating function at the counting point without losing its orbit structure.
\subsection{Sources}
\begin{itemize}
\item[S1] Ulam AI and the UnsolvedMath Contributors, UnsolvedMath: A Curated Collection of Open Mathematics Problems, record 30001568 (OWR-4426-005). Catalogue curation and metadata: CC BY 4.0. The supplied source question is retained; definitions and the English summary are editorial additions.
\url{https://huggingface.co/datasets/ulamai/UnsolvedMath}.
\item[S2] Mini-Workshop: Exploiting Symmetry in Optimization (2011), Oberwolfach Reports, DOI 10.4171/owr/2010/38, pp. 2245--2281
\url{https://doi.org/10.4171/owr/2010/38}.
\textit{Underlying Oberwolfach Reports material retains its source-specific licence, including CC BY-SA 4.0 where applicable; see THIRD\_PARTY\_NOTICES.md.}
\end{itemize}
\paragraph{Preparation and review.}
Original question attributed to the cited source and UnsolvedMath contributors. Codex (GPT-6; exact serving revision not exposed) assisted with definitions, English summaries, source inspection and formatting. No new human mathematical authorship or independent proof verification is asserted. Canonical ID 30001568 and source code OWR-4426-005 checked against the complete supplied registry and existing repository statements; known flagged repetitions and obvious identical conjectures excluded. Full mathematical-equivalence/literature review has not been performed. Statement preparation and final review followed contribution prompts C and D; automated checks verify structure and identity, not open status.
\label{end:30001568}

\section{Literature and context check}
The primary report, printed pages 2267--2268, describes short rational generating functions and explicitly identifies generic evaluation directions as a symmetry difficulty. Its counting application is for a polytope. The attempt below assumes a bounded lattice-point set and an explicitly supplied finite lattice-automorphism group. It preserves orbit blocks by finite averaging, but obtains no savings that avoid expansion over the group.
\section{Research attempt: Orbit-block constant terms and the invariant-direction obstruction}

\paragraph{An exact finite evaluation route}
Suppose a Laurent polynomial $F(z)=\sum_{a\in P\cap\mathbb Z^n}z^a$ for bounded $P$ is given as a finite sum of terms $\varepsilon z^a/\prod_{j=1}^r(1-z^{b_j})$, with nonzero integer $b_j$. Choose an integer vector $v$ with $v\cdot b_j\ne0$ for every denominator vector of every term used. Such a vector can be found by enumerating integer vectors: finitely many proper rational hyperplanes cannot cover all of $\mathbb Z^n$. Explicitly $v=(1,t,\ldots,t^{n-1})$ works for some integer $t$, because each nonzero $b$ yields a nonzero polynomial in $t$ with finitely many roots.

Substitute $z_i=e^{v_i\tau}$. Each rational term becomes a Laurent series with pole order at most $r$. Its constant term can be computed with finitely many rational operations: factor $1-e^{c\tau}=-c\tau(1+c\tau/2+\cdots)$, invert the analytic factor to the required finite degree by recursive coefficient comparison, and multiply with the Taylor series of $e^{(v\cdot a)\tau}$. Terms beyond degree $r$ cannot affect the constant term after the pole of order at most $r$ is extracted. By linearity of coefficient extraction,
\[
 |P\cap\mathbb Z^n|=F(1)=\sum_i\operatorname{CT}_{\tau=0}F_i(e^{v\tau}).
 \tag{1}
\]
The equality holds because the full Laurent polynomial is analytic at $\tau=0$ with value equal to its number of monomials counted with multiplicity one. Individual terms may have poles; their constant terms alone need not be counts.

\paragraph{Keeping a group orbit together}
Let a finite $G\leq\operatorname{GL}_n(\mathbb Z)$ act on exponents. Suppose an orbit block consists of the distinct transforms of a rational term $R$, whose stabilizer as a rational function is $H\leq G$. Choose $v$ generic for all transformed denominators. Its contribution to (1) can be evaluated as
\[
 \frac1{|H|}\sum_{g\in G}\operatorname{CT}_{\tau=0}R(e^{(g^Tv)\tau}).
 \tag{2}
\]
Indeed a transformed exponent $ga$ satisfies $v\cdot ga=(g^Tv)\cdot a$, and each distinct transform is counted exactly $|H|$ times. The entire block value (2) is unchanged if $v$ is replaced by $h^Tv$ for $h\in G$, since this merely permutes the group sum. Thus this finite algorithm retains group-invariant orbit blocks. It can require $|G|$ evaluations per representative and does not solve the requested efficient orbitwise-evaluation challenge. If a block still has an uncanceled pole, its extracted constant term can depend on the generic direction beyond its finite group orbit; only the total in (1) is intrinsically the count.

\paragraph{Why an invariant direction may not exist}
For the segment joining $(1,0)$ and $(0,1)$, the swap group interchanges its two integer points. A rational decomposition is
\[
 F(z_1,z_2)=\frac{z_1}{1-z_2/z_1}+\frac{z_2}{1-z_1/z_2}=z_1+z_2.
\]
Every swap-invariant direction has $v_1=v_2$ and is orthogonal to the denominator exponent $(1,-1)$, so it makes both rational summands undefined identically along the proposed path. A generic $v_1\ne v_2$ works. Since
$1/(1-e^{c\tau})=-1/(c\tau)+1/2+O(\tau)$,
the two constant terms are
\[
 \tfrac12-\frac{v_1}{v_2-v_1},\qquad
 \tfrac12+\frac{v_2}{v_2-v_1},
\]
whose sum is two. This exhibits both the failure of an invariant generic direction and cancellation across an orbit.

Boundedness is essential to interpreting the output as a count. For the nonnegative integer ray, $F(z)=1/(1-z)$ has infinitely many lattice points, whereas this constant-term rule gives $1/2$ along $z=e^\tau$. That number is a regularization, not the cardinality. No regularization convention for general unbounded polyhedra is supplied by the original question.

\section{Verification and adversarial audit}
Genericity is checked for all transformed denominators, not only representatives. Stabilizer division counts distinct rational transforms exactly. The two-point example verifies the Laurent signs and shows that demanding a group-fixed direction can destroy transversality. The unbounded ray separates an algebraic constant term from a finite counting interpretation.
\paragraph{Audit corrections.} Kept pole cancellation, direction dependence of individual blocks, stabilizer multiplicity and bounded-count interpretation explicit.
\paragraph{External computation.} None; symbolic arguments and boundary checks only.
\section{Scope and remaining gaps}
A finite exact orbit-block procedure is proved for bounded input, but it expands over the group and provides no orbit-compression efficiency theorem. Avoiding this expansion while preserving the useful symmetry, and selecting any intended unbounded regularization, remain unresolved. The source's generic-direction method is context rather than claimed novelty.
\paragraph{Integration review.} Codex (GPT-6; exact serving revision not exposed) reviewed the written arguments and expanded layout and source-reference wrappers. Repository metadata and the production web renderer were checked using Python 3.13 and JavaScript with MathJax 3.2.2. This model-assisted review is not independent human proof verification; no PDF compilation or formal proof checking is asserted.

\section{Completion Estimate}
\noindent\textbf{COMPLETION: 0\%}
The calculations below the preserved source statement investigate a concrete route,
but establish no substantiated new progress toward the entire original target.
Known special cases and failed sufficient conditions are not credited as a solution.
\section{References}
The original source references and catalogue attribution are preserved in the
Sources section above. The following URLs were actually consulted for this attempt;
retention of an additional original reference does not claim fresh consultation of it.
\begin{itemize}
\item \url{https://ems.press/content/serial-article-files/46297} (consulted 2026-10-11). 
\end{itemize}
\section{Model and contributor disclosure}
The parent Codex GPT-6 model has reviewed the written mathematical arguments at their declared scopes. This is model-assisted review, not independent human verification.
OpenAI gpt-6-astra, requested reasoning configuration ultra, generated this research
attempt on 2026-10-11. The exact serving revision is not exposed. Codex GPT-6
performed orchestration and repository-format review; its exact serving revision
is also not exposed. The mathematical route, boundary checks and final adversarial
self-audit followed contribution prompts A then B in
\texttt{docs/contribution\_prompts.md}, after reading \texttt{CONTRIBUTING.md}.
Python 3.14.2 assembled the standalone documents; any mathematical computation is
specified in Verification. No theorem prover was used. No human mathematical
authorship, substantial human proof correction, or independent human proof checking
is asserted. Model-assisted review, source retrieval, and structural checks do not
establish correctness or novelty. Full private reasoning transcripts are unavailable
as reproducibility artifacts; the complete arguments, inputs and outcomes actually
used are stated here or in the companion batch audit.
\section{Final status}
Unresolved as a new research contribution. This is a concrete route investigation,
not a claim of new complete resolution or a current-status certification. A known
answer reproduced for context does not become a new contribution. The catalogue's
source-status label is preserved; this attempt status does not assert that a
source-labelled solved question is open.
\end{document}
